\section{Selezione dei Modelli}

La scelta dei modelli adottati è stata effettuata considerando diversi aspetti utili al fine di condurre uno studio adeguato e corretto:

Vengono presentati sei modelli appartenenti a tre laboratori di ricerca diversi: Alibaba Cloud, Meta e DeepSeek AI. 
I modelli sono tutti di medie-piccole dimensioni essendo nell'ordine dei 7 miliardi di parametri. Sono stati scelti per il loro largo impiego nel mondo aziendale, per la loro diversità architetturale e per essere adatti a prestarsi a diversi utilizzi come: assistente IA, sviluppo software o matematica complessa e logica.

La scelta di modelli di questa grandezza è dovuta principalmente a ragioni computazionali: la natura della ricerca, in particolar modo della Mechanistic Interpretability, richiede l'estrazione e la manipolazione sistematica di tensori ad alta dimensionalità da ogni layer. L'impiego di questi modelli ha permesso di mantenere la corretta esecuzione garantendo un'elevata capacità di ragionamento semantico con tempistiche adeguate. I modelli selezionati per lo studio sono:\\

\textbf{Famiglia Qwen (Alibaba)}
\begin{itemize}[noitemsep]
    \item Qwen2.5-7B-Instruct
    \item Qwen2.5-Coder-7B-Instruct
\end{itemize}

\textbf{Famiglia Llama (Meta)}
\begin{itemize}[noitemsep]
    \item Llama-3.1-8B-Instruct
    \item CodeLlama-7b-Instruct-hf
\end{itemize}

\textbf{Famiglia DeepSeek (DeepSeek)}
\begin{itemize}[noitemsep]
    \item DeepSeek-R1-Distill-Qwen-7B
    \item deepseek-coder-6.7b-instruct
\end{itemize}

Per poter portare un maggior numero di scenari cinque di questi modelli sono a stampo Instruction-tuned cercando di dare la risposta finale e corretta nel modo più rapido ed efficiente possibile, senza mostrare i calcoli interni a meno di una richiesta esplicita.
Il modello DeepSeek-R1-Distill-Qwen-7B presenta invece una Chain-of-Thought nativa: addestrato con tecniche di reinforcement learning per ragionare obbligatoriamente prima di rispondere generando una catena di pensieri prima di consegnare la risposta finale.

Tutti i modelli presentano un'architettura decoder-only e, al fine di mantenere più eterogeneità possibile, per lo studio sono stati adottati quattro modelli appartenenti allo stato dell'arte odierno, quali Qwen2.5-7B-Instruct, Llama-3.1-8B-Instruct, Qwen2.5-Coder-7B-Instruct e DeepSeek-R1-Distill-Qwen-7B, e due modelli più vecchi quali deepseek-coder-6.7b-instruct e CodeLlama-7b-Instruct-hf così da avere anche un riscontro temporale dello studio e analizzare se le vulnerabilità sono state nascoste o superate.

Per ognuno dei laboratori di ricerca sono stati scelti un modello base e un modello specializzato per il coding: poiché la ricerca si basa sulla sicurezza del codice presentato ad un LLM è necessario osservare, analizzare e capire la natura e il comportamento di due modelli addestrati in maniera differente e come questi rispondano a perturbazioni esterne.

\subsubsection{Famiglia Qwen}

Qwen2.5 è lo standard serie di Qwen Large Language Models. Qwen 2.5 ha portato diversi miglioramenti rispetto ai predecessori: molta più conoscenza e ha migliorato le capacità di Coding e matematiche assieme a miglioramenti in Instruction Following, generazione di testi lunghi oltre 8k tokens, comprendere dati strutturati e generarne di nuovi.
Qwen2.5-7B è un modello Instruction-Tuned con circa 7.61 miliardi di parametri e una Context Window di 128.000 token, ha 28 layer e la capacità di generare 8192 token.
L'architettura adottata è quella di un Transformers con RoPE, SwiGLU, RMSNorm e attention QKV bias \cite{qwen2.5-hf}.
Qwen2.5-coder è la serie di Qwen specifica per il Coding che ha portato diversi miglioramenti nella generazione, nel Reasoning e nella correzione di script.
Come il modello di base è composto da 28 layer e ha 7.61 miliardi di parametri e una Context Windows di 128.000 Token e condivide con esso la stessa architettura Transformer \cite{hui2024qwen2-coder-hf}.

\subsubsection{Famiglia Llama}

Llama 3.1 rappresenta una delle ultime generazioni dei Large Language Models open-weight sviluppati da Meta. Rispetto ai suoi predecessori, introduce significativi miglioramenti nelle capacità di ragionamento multilingue, nella lunghezza del contesto e, fondamentale per questo studio, nell'allineamento alla sicurezza (Safety Alignment) tramite tecniche avanzate di RLHF (Reinforcement Learning from Human Feedback) e filtri di rifiuto più severi.
Llama-3.1-8B-Instruct è un modello generalista Instruction-Tuned che dispone di circa 8 miliardi di parametri. A differenza dell'architettura Qwen, presenta 32 Layer interni. Supporta una Context Window estesa a 128.000 Token ed è basato su un'architettura Transformer standard ottimizzata con Grouped-Query Attention (GQA) per migliorare l'efficienza computazionale \cite{Llama-3.1-8B-Instruct-hf}.
CodeLlama-7b-Instruct-hf è, invece, una release precedente specializzata nella generazione e discussione di codice. Pur essendo un modello Instruction-Tuned, si basa sulla passata architettura di Llama 2. Mantiene una struttura a 32 layer con 7 miliardi di parametri, ma con una Context Window limitata a 100.000 token. L'inclusione di questo modello funge da Baseline storica, essenziale per analizzare come le tecniche di sicurezza si siano evolute nel tempo \cite{CodeLlama-7b-Instruct-hf}.

\subsubsection{Famiglia DeepSeek}

I modelli DeepSeek offrono un punto di osservazione unico per questo studio, includendo sia architetture proprietarie che ibride.
DeepSeek-Coder-6.7b-instruct è un modello specializzato nel codice, addestrato su un vastissimo corpus di repository progettuali (87\% codice, 13\% linguaggio naturale in inglese e cinese). Questo modello instruction-tuned si basa su un'architettura Transformer proprietaria a 32 layer, ottimizzata per finestre di contesto da 16.000 token, con estensioni fino a 32k. Nonostante sia tecnologicamente precedente agli attuali modelli SOTA, la sua elevatissima competenza semantica sul codice lo rende un bersaglio primario per l'analisi meccanicistica \cite{deepseekai2025deepseekr1incentivizingreasoningcapability-hf}.
Il modello DeepSeek-R1-Distill-Qwen-7B introduce, infine, un cambio di paradigma radicale per questo studio. Questo modello non utilizza l'architettura base di DeepSeek, bensì i pesi base di Qwen 2.5, mantenendo quindi 28 layer e 7 miliardi di parametri. Tuttavia, è stato "distillato" utilizzando gli output del modello "ragionatore" DeepSeek-R1 originale. Questo processo ha infuso nel modello un comportamento Chain-of-Thought nativo: di fronte a un problema, il modello è forzato a generare un blocco latente di ragionamento logico prima di produrre la risposta finale \cite{deepseek-coder-6.7b-instruct-hf}.